
\section{Le principe de localité des données}

\begin{frame}{Les programmes vers les données, pas l'inverse!}

Client serveur pour traiter des TOs de données? \alert{Ne marche pas}.

\vfill

\figSlideWithSize{data-locality}{13cm}

\vfill

Il faut placer les traitements au plus près des données.

\end{frame}


\begin{frame}{Pourquoi un \textit{framework}}

Un framework \alert{pilote} un traitement, conduit selon un processus générique (notion d'inversion de contrôle).


\vfill

L'application ``injecte'' des fonctions dans le framework.

% Exemples: Applications MVC, XSLT, Couches ORM, \alert{MapReduce}.
 
\vfill

Le framework applique les fonctions  pendant l'exécution du processus.

\vfill

\begin{redbullet}
  \item Une fonction \textit{map()} à appliquer pendant la phase de Map
  \item Une fonction \textit{reduce()} à appliquer pendant la phase de Reduce.
\end{redbullet}

\begin{remblock}{Rôle d'un \textit{framework} MapReduce}
Prendre en charge la distribution, et gérer les reprises sur panne.
\end{remblock}
\end{frame}


\begin{frame}[fragile]{Comptons les mots: la fonction de Map}

Phase de Map: on extrait les termes, on les compte (localement), on transmet

\begin{minted}{bash}
    function mapTF($id, $contenu)
    {
      // $id: identifiant du document
      // $contenu: contenu textuel du document

      // On boucle sur tous les termes du contenu
      foreach  ($t in $contenu) {
        // Comptons le nb d'occurrences de $t dans $contenu
        $count = nbOcc ($t, $contenu);
        // Emission du terme et de son nombre d'occurrences
        emit ($t, $count);
      }
    }
\end{minted}

NB: rien n'indique le contexte de distribution.

\end{frame}


\begin{frame}[fragile]{Comptons les mots: la fonction de Reduce}

Phase de Reduce: on reçoit, pour chaque terme, les compteurs, on les additionne.

\begin{minted}{bash}
    function reduceTF($t, $compteurs)
    {
      // $t: un terme
      // $compteurs: les nombres d'occurrences, un pour chaque doc.
      $total = 0;
      // Boucles sur les compteurs et calcul du total
      foreach ($c in $compteurs) {
       $total = $total + $c;
      }
      // Et on produit le total  
      return $total;
    }
\end{minted}

Même remarque: rien n'indique le contexte de distribution.

\end{frame}



\newcommand{\map}{\mathop{\mathit{map}}}
\newcommand{\reduce}{\mathop{\mathit{reduce}}}
\newcommand{\mapreduce}{\mathop{\mathit{MapReduce}}}
\newcommand{\shuffle}{\mathop{\mathit{shuffle}}}
\newcommand{\combiner}{\mathop{\mathit{combiner}}}

\begin{frame}{L'exécution par le \textit{framework}}

\begin{tabular}{ll}
URL&Document\\ \hline
$u_1$&the jaguar is a new world mammal of the felidae family.\\
$u_2$&for jaguar, atari was keen to use a 68k family device.\\
$u_3$&mac os x jaguar is available at a price of us \$199 for
apple's\\& new ``family pack''.\\
$u_4$&one such ruling family to incorporate the jaguar into their name\\&is
jaguar paw.\\
$u_5$&it is a big cat.\\
\hline
\end{tabular}

\end{frame}

\begin{frame}{Déroulons le processus}

\begin{columns}[t]
\column{.33\linewidth}
\centering\begin{tabular}{ll}
\hline
\structure{term}&\structure{count}\\
\hline
jaguar&1\\
mammal&1\\
family&1\\
jaguar&1\\
available&1\\
jaguar&1\\
family&1\\
family&1\\
jaguar&2\\
\dots\\
\hline
\end{tabular}
Sortie du $\map$\\
Entrée du $\shuffle$
\pause
\column{.33\linewidth}
\centering\begin{tabular}{ll}
\hline
\structure{term}&\structure{count}\\
\hline
jaguar&1,1,1,2\\
mammal&1\\
family&1,1,1\\
available&1\\
\dots\\
\hline
\end{tabular}
Sortie du $\shuffle$ \\
Entrée du $\reduce$ 
\pause
\column{.33\linewidth}
\centering\begin{tabular}{ll}
\hline
\structure{term}&\structure{count}\\
\hline
jaguar&5\\
mammal&1\\
family&3\\
available&1\\
\dots\\
\hline
\end{tabular}
Résultat final
\end{columns}

\end{frame}


\section{Exécution distribuée d'un traitement \mapreduce/}

\begin{frame}{La grande vision}

\figSlideWithSize{mr-execution}{11cm}

\end{frame}


\begin{frame}
\frametitle{Quelques calculs}

Supposons une collection avec 1 million de documents; chaque document contient 100 termes
en moyenne, et sa taille est (en moyenne) de 1 KO.

\vfill

Le partitionnement découpe en fragments de 64 MO: chaque fragment 
contient  64\ 000 documents. Il y a à peu près $M= \lceil 1{,}000{,}000/64{,}000
\rceil \approx 16{,}000$ fragments.

\vfill

Si j'ai 16 machines, chacune traitera 1\ 000 fragments par 1\ 000 tâches.

\vfill

Chaque tâche (Mapper) produit 6\, 400\ 000 paires (terme, compteur)

\vfill

Si j'ai 10 machines pour la phase de Reduce: chaque Reducer $R_i$ reçoit 640\ 000 paires
de chaque Mapper, pour tous les termes $t$ tels que  $\mathit{hash}(t)=i$ ($1\leq i\leq {10}$)

\end{frame}



\begin{frame}{Un tout petit exemple}

Quatre documents qui parlent de loups, de moutons, de bergerie.
\vfill


\figSlideWithSize{mr-execution-ex}{11cm}

\end{frame}


\section{La reprise sur panne}


\begin{frame}
\frametitle{Gestion des pannes}

Un traitement MapReduce peut durer des jours, sur des centaines de machines:
la probabilité de panne est très grande.

\vfill

Ré-exécuter à chaque panne $\Rightarrow$ : on ne finira jamais
\vfill

Dans Hadoop, le Maître (\textit{JobTracker}) surveille l'ensemble du processus.

\medskip

\begin{redbullet}
  \item Panne d'un Reducer: on  reprend le résultat des Mappers, \alert{stocké sur disque}.

  \item Panne d'un Mapper: on recommence la tâche sur le fragment (\alert{ou sur un réplica})
  \item Si le Maître tombe en panne: on recommence tout, (probabilité très faible).

\end{redbullet}

\begin{remblock}{Essentiel}
Tout repose sur une \alert{sérialisation} (écriture sur disque) aux différentes
étapes. \alert{Robuste mais très lent}.
\end{remblock}
\end{frame}
